\documentclass[10pt,a4paper]{amsart}
\usepackage[margin=19mm]{geometry}
\usepackage{amsmath,amssymb,mathtools}
\usepackage[scr=rsfs,cal=euler]{mathalfa}
\usepackage{xfrac}
\usepackage[unicode,hidelinks,bookmarks=false]{hyperref}
\newcommand{\ie}{i.e.\ }
\newcommand{\cf}{cf.\ }
\newcommand{\resp}{respectively\ }
\newcommand{\Subsection}{\subsection}
\providecommand{\nobreakdash}{\nobreak\hbox{-}}
\setlength{\emergencystretch}{2em}
\multlinegap0pt
\usepackage{amssymb}
\usepackage{tikz}
\usepackage{xcolor}
\newtheorem{theorem}{Theorem}
\newtheorem{lemma}{Lemma}
\title{Confirmation of the Daykin-Frankl Conjecture}
\author{Kada Williams}
\date{Source: arXiv:2609.03087v1 (CC BY 4.0)}
\begin{document}
\maketitle

\noindent\emph{Source and licence.} Confirmation of the Daykin-Frankl Conjecture. Version arXiv:2609.03087v1 (CC BY 4.0), distributed by arXiv under the Creative Commons Attribution 4.0 International licence, http://creativecommons.org/licenses/by/4.0/, as recorded in the arXiv licence metadata for this exact version and shown on the abstract page https://arxiv.org/abs/2609.03087v1. Full licence notice: \texttt{licenses/arxiv-2609.03087-CC-BY-4.0.txt}.\par\smallskip

\noindent\textit{Editorial scope.} Counted unit: the article's lemma key (source0.tex lines 74--77). The article's complete original proof of it is delivered with it (source0.tex lines 79--85, one \texttt{proof} environment of 7 lines). No mathematics is written, repaired or re-derived: the statement and the proof are the article's own lines, in the article's own notation, and no retained line is substituted or re-flowed.  Retained uncounted as the source-local context the proof needs: lines 67--70.  Nothing was omitted from any retained span, so the package carries no omission record.  No retained line contains a citation, so the wrapper carries no bibliography and no bibliography entry.  No retained line contains a figure environment, an image-inclusion command or drawing code, so no image is delivered and the package declares no asset dependency.  Editorial formatting only, in addition: an AMS \texttt{amsart} preamble that restates the article's own declarations and macros used by the retained lines, namely the environments \texttt{theorem}, \texttt{lemma} on separate counters, together with the standard packages those lines need.  The retained environments print in this document as Theorem 1, Lemma 1 in order of appearance, whereas the article prints the same items with the numbering of its own full document.  The complete native source of the article as distributed in the arXiv e-print for this exact version is bundled unchanged as \texttt{sources/209262/source0.tex}.
\begin{theorem} \label{main}
    Let $n,k\ge 0$ be integers. Suppose $P\subseteq Q_n$ is a convex subset of $Q_n$, ordered by the subset relation. Then
    $$w(P\times Q_k)\ge w(Q_{n+k})|P|2^{-n}.$$
\end{theorem}

\begin{lemma} \label{key}
    Let $R$ be a poset and $P\subseteq R\times Q_1$ be a convex subset. We identify $Q_1$ as $\{0,1\}$ and let $P=P_0\times \{0\}\cup P_1\times \{1\}$, where $P_0,P_1\subseteq R$. Then
    $$w(P)\ge \frac{w(P_0\times Q_1)+w(P_1\times Q_1)}{2}.$$
\end{lemma}

\begin{proof}[Proof of Theorem \ref{main} via Lemma \ref{key}]
    We induct on $n$. If $n=0$, $Q_0=\{\emptyset\}$, so $P=\emptyset$ or $P=Q_0$, with $w(P\times Q_k)=0$ or $=w(Q_k)$, respectively, as claimed. Let us now prove the result for $n+1$, given that it holds for $n$.

    Let $P\subseteq Q_{n+1}\cong Q_n\times Q_1$ be convex and $P=P_0\times \{0\}\cup P_1\times\{1\}$. Then also $P\times Q_k\subseteq Q_{n+1}\times Q_k\cong Q_{n+k}\times Q_1$ is convex. By Lemma \ref{key}, with $R=Q_{n+k}$,
    $$w(P\times Q_k)\ge \frac{w(P_0\times Q_{k+1})+w(P_1\times Q_{k+1})}{2}.$$
    Since $P_0,P_1\subseteq Q_n$ are convex, the result follows by induction.    
\end{proof}

\end{document}
